\section*{第 \ref{ch:b3:hilbert} 章——希尔伯特空间}

\begin{solution}{exo:b3:hilbert:1}
(a) 实数情形：展开 $\norm{x \pm y}^2 = \norm x^2 \pm 2\langle
x,y\rangle + \norm y^2$，相减即得。复情形（\omterm{def:b3:hilbert:inner}{内积}
关于第二变元线性）：仿上式展开，
\[
\langle x, y\rangle = \frac14\sum_{k=0}^{3}
\iu^k\,\bigl\|\iu^kx + y\bigr\|^2,
\]
每一项贡献 $\iu^k\cdot2\operatorname{Re}\bigl(
(-\iu)^k\langle x,y\rangle\bigr)$，其和为 $4\langle
x,y\rangle$（逐一检验 $k$ 的四个取值；$\sum\iu^k
(\norm x^2 + \norm y^2) = 0$）。

(b) $L^1$ 情形：$f = \mathbf 1_{\intcc0{1/2}}$，$g = \mathbf
1_{\intcc{1/2}1}$：$\norm{f\pm g}_1^2 = 1$（两式各自），和为 $2$；
$2\norm f_1^2 + 2\norm g_1^2 = 1 \neq 2$。上确界范数：$f =
\mathbf 1$，$g(t) = t$ 定义于 $\intcc01$：$\norm{f + g}_\infty^2 +
\norm{f-g}_\infty^2 = 4 + 1 = 5 \neq 4 = 2 + 2$。由于不满足
平行四边形法则，这些范数都不由任何\omterm{def:b3:hilbert:inner}{内积}诱导
（\omterm{def:b3:hilbert:inner}{内积}必满足此式，直接展开即可验证）。
\end{solution}

\begin{solution}{exo:b3:hilbert:2}
(a) $p(f) = \bigl(\int_0^1f\bigr)\mathbf 1$：事实上 $f - \int f
\perp$ 与常数正交（$\int(f - \int f)c = 0$）。在均方意义下，
$f$ 的最佳常数逼近即其\emph{平均值}——
这是条件期望的第一个实例
（见\cref{ch:b3:probability}）。

(b) $p(f) = f\,\mathbf 1_{\intcc{1/2}1}$：其差
$f\mathbf 1_{\intcc0{1/2}}$ 与每个 $g$
（在 $\intcc0{1/2}$ 上为零）正交。

(c) $d^2 = \bigl\|x - \tfrac12\bigr\|_2^2 = \int_0^1(x -
\tfrac12)^2\dd x = \tfrac1{12}$：$d = \frac1{2\sqrt3}$。
\end{solution}

\begin{solution}{exo:b3:hilbert:3}
(a) 该等价性即\cref{thm:b3:hilbert:decomposition}
（$\bar F = (F^\perp)^\perp$，且 $\bar F = H \iff F^\perp =
\{0\}$）。例：由有限数列组成的空间 $F$ 在
$\ell^2$ 中稠密（截断），且为真子空间：$F^\perp = \{0\}$，但 $F
\neq \ell^2$——对于非闭子空间，$H = F \oplus
F^\perp$ 显然不成立（$F \oplus \{0\} \neq H$）。

(b) $\abs{\langle x_n, y_n\rangle - \langle x, y\rangle} \leq
\abs{\langle x_n - x, y_n\rangle} + \abs{\langle x, y_n -
y\rangle} \leq \norm{x_n - x}\sup_n\norm{y_n} + \norm
x\,\norm{y_n - y} \to 0$（收敛数列有界）。
用于 \cref{thm:b3:hilbert:parseval}(1) 中，说明 $x - \lim
S_N \perp e_k$，并用于 \cref{thm:b3:hilbert:decomposition}
中，说明 $F^\perp = \bar F^{\,\perp}$。
\end{solution}

\begin{solution}{exo:b3:hilbert:4}
$e_0 = \frac1{\sqrt2}$。其次，$x \perp \mathbf 1$ 已成立
（$\int_{-1}^1x = 0$），且 $\int_{-1}^1x^2 = \frac23$：$e_1 =
\sqrt{\tfrac32}\,x$。再取 $x^2 - \langle e_0, x^2\rangle e_0 =
x^2 - \frac13$（由奇偶性知 $\perp e_1$），并有
\[
\int_{-1}^1\Bigl(x^2 - \frac13\Bigr)^2\dd x = \frac25 -
\frac49 + \frac29 = \frac{8}{45}:
\qquad e_2 = \sqrt{\tfrac{45}8}\,\Bigl(x^2 - \frac13\Bigr).
\]
比较：$P_0 = 1$，$P_1 = x$，$P_2 = \frac{3x^2 - 1}2$，则
$\sqrt{n + \tfrac12}\,P_n$ 给出 $\frac1{\sqrt2}$，
$\sqrt{\frac32}x$，$\sqrt{\frac52}\,\frac{3x^2-1}2 =
\sqrt{\frac{45}8}\bigl(x^2 - \frac13\bigr)$，恰为 $e_0, e_1,
e_2$。
\end{solution}

\begin{solution}{exo:b3:hilbert:5}
对 $f(t) = t$：$c_0 = 0$，由分部积分得 $c_n =
\frac{\iu(-1)^n}{n}$，当 $n \neq 0$ 时，$\abs{c_n}^2 =
\frac1{n^2}$。由\omterm{thm:b3:hilbert:parseval}{帕塞瓦尔恒等式}：
\[
\frac1{2\pi}\int_{-\pi}^\pi t^2\dd t = \frac{\pi^2}3
= \sum_{n\neq0}\frac1{n^2} = 2\sum_{n\geq1}\frac1{n^2}
\ \Longrightarrow\ \sum_{n\geq1}\frac1{n^2} = \frac{\pi^2}6 .
\]
对 $f(t) = t^2$：$c_0 = \frac{\pi^2}3$，$c_n =
\frac{2(-1)^n}{n^2}$（$n \ne 0$）。由\omterm{thm:b3:hilbert:parseval}{帕塞瓦尔恒等式}：
\[
\frac1{2\pi}\int_{-\pi}^{\pi}t^4\dd t = \frac{\pi^4}5
= \frac{\pi^4}9 + \sum_{n\neq0}\frac4{n^4}
\ \Longrightarrow\
\sum_{n\geq1}\frac1{n^4} = \frac18\Bigl(\frac{\pi^4}5 -
\frac{\pi^4}9\Bigr) = \frac{\pi^4}{90} .
\]
无需分段 $\mathcal C^1$ 之类的附加条件：
由\cref{thm:b3:hilbert:fourier}即可覆盖每个 $L^2$ 函数。
\end{solution}

\begin{solution}{exo:b3:hilbert:6}
(a) $\varphi(f) = \int_0^{1/2}f = \langle\mathbf
1_{\intcc0{1/2}},\ f\rangle$：代表元为 $a =
\mathbf 1_{\intcc0{1/2}}$，且 $\norm\varphi = \norm a_2 =
\frac1{\sqrt2}$（见\cref{thm:b3:hilbert:riesz}）。

(b) 取帐篷函数 $f_n$，其峰值 $1$ 位于
$\frac12$ 处，支集宽为 $\frac2n$：$f_n(\tfrac12) = 1$
而 $\norm{f_n}_2^2 \leq \frac2n \to 0$，故不存在常数 $C$
使得 $\abs{f(\frac12)} \leq C\norm f_2$。点求值
在 $L^2$ 中没有意义——元素是模去
零测集的等价类，而本题的计算正是其定量上的原因。
\end{solution}

\begin{solution}{exo:b3:hilbert:7}
(a) 记 $M = \sup_k\norm{x_k}$。标量数列
$(\langle e_n, x_k\rangle)_k$ 均以 $M$ 为界：由对角线
抽取法可得子列 $x_{k_j}$，使得 $\langle e_n, x_{k_j}\rangle
\to \gamma_n$（对每个 $n$）。对每个 $N$，$\sum_{n\leq
N}\abs{\gamma_n}^2 = \lim_j\sum_{n\leq N}\abs{\langle e_n,
x_{k_j}\rangle}^2 \leq M^2$（由贝塞尔不等式），故 $(\gamma_n) \in
\ell^2$，且 $x = \sum_n\gamma_ne_n \in H$
（见\cref{thm:b3:hilbert:parseval}(3)）。对 $y \in H$：
\[
\abs{\langle y, x_{k_j} - x\rangle}
\leq \Bigl|\sum_{n\leq N}\overline{c_n(y)}\bigl(\langle e_n,
x_{k_j}\rangle - \gamma_n\bigr)\Bigr|
+ 2M\Bigl(\sum_{n>N}\abs{c_n(y)}^2\Bigr)^{1/2},
\]
利用展开式 $\langle y, z\rangle =
\sum\overline{c_n(y)}c_n(z)$，并对尾项使用柯西–施瓦茨不等式；
先取 $N$，再取 $j$：弱收敛于 $x$。

(b) $\langle y, e_n\rangle = c_n(y) \to 0$ 对每个 $y$
（$\ell^2$ 尾部）：$e_n \rightharpoonup 0$，但 $\norm{e_n} =
1$：范数不是弱连续的。在 (a) 中：$\norm x^2 =
\sum\abs{\gamma_n}^2 \leq \liminf_j\norm{x_{k_j}}^2$（有限
截断及再次用贝塞尔不等式）：弱极限只会使范数减小。
\end{solution}

\begin{solution}{exo:b3:hilbert:8}
固定 $y$，$x \mapsto \langle y, Tx\rangle$ 是一个
连续线性泛函；由里斯表示定理，存在唯一的 $T^*y$ 使
$\langle y, Tx\rangle = \langle T^*y, x\rangle$ 对所有 $x$ 成立
——取共轭得 $\langle Tx, y\rangle = \langle x,
T^*y\rangle$。由唯一性知 $T^*$ 线性；
\[
\norm{T^*y} = \sup_{\norm x = 1}\abs{\langle T^*y, x\rangle}
= \sup_{\norm x=1}\abs{\langle y, Tx\rangle}
\leq \vertiii T\,\norm y,
\]
于是 $\vertiii{T^*} \leq \vertiii T$，由 $T^{**} = T$
得等号。移位算子：$\langle Sx, y\rangle = \sum_{n\geq1}
x_n\bar y_{n+1} = \langle x, S^*y\rangle$，其中 $(S^*y)_n =
y_{n+1}$：后向移位算子。核--像：$T^*y = 0$ 当且仅当
$\langle x, T^*y\rangle = 0$ 对所有 $x$ 成立，当且仅当 $\langle Tx,
y\rangle = 0$ 对所有 $x$ 成立，当且仅当 $y \perp \operatorname{im}T$：
$\ker T^* = (\operatorname{im}T)^\perp$；两边取 $\perp$，并
利用\cref{thm:b3:hilbert:decomposition}，
$\overline{\operatorname{im}T} = (\ker T^*)^\perp$。
\end{solution}

\begin{solution}{exo:b3:hilbert:9}
若 $a(u, \cdot) = \varphi$，则对任意 $w$，
\[
J(u + w) - J(u) = a(u, w) - \varphi(w) + \tfrac12a(w,w)
= \tfrac12a(w,w) \geq \tfrac\alpha2\norm w^2,
\]
当 $w \neq 0$ 时严格为正，故 $u$ 为唯一的极小点。
反之，在极小点处，函数 $t \mapsto J(u + tw)$
（关于 $t$ 的二次多项式）在
$0$ 处导数为零：$a(u, w) - \varphi(w) = 0$ 对每个 $w$ 成立。
投影的重新推导：对闭子空间 $F$，在\omterm{def:b3:hilbert:inner}{希尔伯特空间} $F$ 上应用拉克斯–米尔格拉姆定理，
取 $a(u,v) = \langle u, v\rangle$
（$M = \alpha = 1$），且 $\varphi(v) = \langle x, v\rangle$：存在
唯一的 $p \in F$ 使得 $\langle p, v\rangle = \langle x,
v\rangle$ 对所有 $v \in F$ 成立，即 $x - p \perp F$——
由对称情形，$p$ 使 $\frac12\norm v^2 - \langle
x, v\rangle = \frac12\norm{v - x}^2 - \frac12\norm x^2$ 在
$F$ 上取极小，即为投影。
\end{solution}

\begin{solution}{exo:b3:hilbert:10}
规范化：$\int h_n^2 = 2^j\cdot 2^{-j} = 1$。
正交性：两个不同的哈尔函数要么
支集的\omterm{def:b3:topology:interior}{内部}互不相交（乘积几乎处处为零），或者
较细者的支集含于一个半区间，在该半区间上较粗者
为常数——于是乘积的积分为
该常数乘 $\int h_{\text{较细}} = 0$；对 $h_0
= \mathbf 1$，同样 $\int h_n = 0$。完备性：
$\{h_0, \dots, h_{2^J-1}\}$ 的张成空间是
步长为 $2^{-J}$ 的二进网格上的阶梯函数；两空间维数均为
$2^J$，且哈尔函数线性无关（即标准正交）：
故张成空间恰为\emph{全部}此类阶梯函数。二进
阶梯函数在 $L^2(\intcc01)$ 中稠密：\omterm{def:b3:lebesgue:simple}{简单函数}
在\cref{thm:b3:lp:density}(1)中稠密，可测集
可由有限个区间之并逼近
（见\cref{exo:b3:measure:7}），区间则
用二进区间逼近（端点移动量 $\leq 2^{-J}$）。由
\cref{thm:b3:hilbert:parseval}，哈尔系是\omterm{def:b3:hilbert:onb}{希尔伯特
基}。
\end{solution}

\begin{solution}{exo:b3:hilbert:11}
(i) $\Rightarrow$ (ii)：对正交投影，
$\langle Px, y\rangle = \langle Px, Py\rangle = \langle x,
Py\rangle$（代入分解 $x = Px + (x - Px)$
等，并消去交叉项）。(ii) $\Rightarrow$ (iii)：
$\norm{Px}^2 = \langle P^2x, x\rangle = \langle Px, x\rangle
\leq \norm{Px}\norm x$，故 $\vertiii P \leq 1$，且 $Pu = u$
在非零像上：$= 1$。(iii) $\Rightarrow$ (i)：$H =
\operatorname{im}P \oplus \ker P$（代数上，由 $P^2
= P$）；取 $u = Pu \in \operatorname{im}P$，$v \in \ker P$，
$t \in \R$：$\norm{P(u + tv)}^2 = \norm u^2$ 必须 $\leq
\norm{u + tv}^2 = \norm u^2 + 2t\operatorname{Re}\langle u,
v\rangle + t^2\norm v^2$ 对每个 $t$ 成立，迫使
$\operatorname{Re}\langle u, v\rangle = 0$（比较
$t \to 0^\pm$ 时的一次项）；将 $v$ 换成 $\iu v$，
亦消去虚部：$\operatorname{im}P \perp
\ker P$，这恰是投影的正交性。
例：$P(x, y) = (x + y, 0)$ 在 $\R^2$ 上：$P^2 = P$，
像为 $x$ 轴，核为直线 $y = -x$，且
$\vertiii P = \sup\frac{\abs{x+y}}{\norm{(x,y)}} = \sqrt2$
（在 $(1,1)/\sqrt2$ 处取得）：斜投影的范数
$> 1$。（补充：(ii) 也直接推出 (i)：
$\ker P = \ker P^* = (\operatorname{im}P)^\perp$，由
\cref{exo:b3:hilbert:8}）。
\end{solution}

\begin{solution}{exo:b3:hilbert:12}
(a) 对酉算子 $U$：$\norm{Ux - x}^2 = 2\norm x^2 -
2\operatorname{Re}\langle Ux, x\rangle$ 及 $\norm{U^*x -
x}^2 = 2\norm x^2 - 2\operatorname{Re}\langle x, Ux\rangle$：
二者同时为零，故 $\ker(U - I) = \ker(U^* - I)$。
再利用 $\ker T^* = (\operatorname{im}T)^\perp$
（见\cref{exo:b3:hilbert:8}），取 $T = U - I$，$T^* = U^* -
I$：
\[
\overline{\operatorname{im}(U - I)} = \bigl(\ker(U^* -
I)\bigr)^\perp = F^\perp .
\]

(b) 在 $F$ 上：$U^kx = x$，故 $A_nx = x$。对 $x = (U - I)y$：
$A_nx = \frac1n(U^ny - y)$，其范数 $\leq \frac2n\norm y \to
0$。对\omterm{def:b3:topology:interior}{闭包}中的 $x$：给定 $\varepsilon$，取 $x' =
(U - I)y$，使 $\norm{x - x'} < \varepsilon$；由于
$\vertiii{A_n} \leq \frac1n\sum\vertiii{U^k} = 1$，
$\norm{A_nx} \leq \norm{A_n(x - x')} + \norm{A_nx'} \leq
\varepsilon + o(1)$。

（c）将 $x = Px + (x - Px)$ 分解，其中 $Px \in F$ 且 $x -
Px \in F^\perp = \overline{\operatorname{im}(U - I)}$（由
(a) 部分）：$A_nx = Px + A_n(x - Px) \to Px + 0$。

（d）在傅里叶基 $e_m(x) = \eu^{2\iu\pi mx}$ 中：$Ue_m
= \eu^{2\iu\pi m\alpha}e_m$，故 $Ue_m = e_m$ 当且仅当 $m\alpha
\in \Z$ 当且仅当 $m = 0$（$\alpha$ 为无理数）：$F = \C\mathbf 1$
且 $Pf = \langle\mathbf 1, f\rangle\mathbf 1 = \int_0^1f$。
定理即为 $\frac1n\sum_{k<n}f(\cdot + k\alpha) \to
\int_0^1f$ 在 $L^2(\R/\Z)$ 中：无理旋转的\omterm{def:b3:groups:action}{轨道}平均
呈等分布——此为 $L^2$ 意义下的
外尔等分布定理的影子，由纯粹的
希尔伯特几何即可得到。
\end{solution}

\begin{solution}{pb:b3:hilbert:1}
\textbf{1.} 格拉姆--施密特正交化保证
$\operatorname{Vect}(p_0, \dots, p_n) =
\operatorname{Vect}(1, \dots, t^n) = \R_n[t]$ 且 $p_n \perp
p_k$（$k < n$），从而 $p_n \perp \R_{n-1}[t]$。$p_k$ 的
次数严格递增，线性无关，构成基。

\textbf{2.} $t\,p_n$ 是次数为 $n + 1$ 的首一多项式：展开
$t\,p_n = p_{n+1} + \sum_{k\leq n}c_kp_k$，其中 $c_k =
\langle p_k, tp_n\rangle/\norm{p_k}^2$。对 $k \leq n - 2$：
$\langle p_k, tp_n\rangle = \langle tp_k, p_n\rangle = 0$
（次数 $k + 1 < n$）。于是 $tp_n = p_{n+1} + a_np_n +
b_np_{n-1}$，即所述递推式，其中
\[
b_n = \frac{\langle p_{n-1}, tp_n\rangle}{\norm{p_{n-1}}^2}
= \frac{\langle tp_{n-1}, p_n\rangle}{\norm{p_{n-1}}^2}
= \frac{\langle p_n + (\text{低次项}),\
p_n\rangle}{\norm{p_{n-1}}^2}
= \frac{\norm{p_n}^2}{\norm{p_{n-1}}^2} > 0 .
\]

\textbf{3.} 设 $t_1 < \dots < t_m$ 为
$I$ \omterm{def:b3:topology:interior}{内部}使 $p_n$ 变号的点，并令 $q = \prod_{i\leq m}(t -
t_i)$（规定 $q = 1$，当 $m = 0$ 时）。于是 $p_nq$ 在
$I$ 上符号恒定，且几乎处处不为零：$\int_Ip_nq\,w \neq 0$。若 $m <
n$，则与 $p_n \perp \R_{n-1}[t]$ 矛盾。于是 $m = n$：
$p_n$ 恰有 $n$ 个互异的\omterm{def:b3:topology:interior}{内部}根（它至多有 $n$
个根）。

\textbf{4.} $(t^2 - 1)^n$ 的次数为 $2n$；$n$ 次求导
后次数仍为 $n$，首项系数为
$\frac{(2n)(2n-1)\cdots(n+1)}{2^nn!} =
\frac{(2n)!}{2^n(n!)^2}$。对 $\deg Q < n$，分部积分
$n$ 次：所有边界项都含阶数
$< n$ 的 $(t^2-1)^n$ 导数，它在 $\pm1$ 处为零（
$n$ 重根）。经 $n$ 步后，被积函数含 $Q^{(n)} = 0$。

\textbf{5.} 记 $u = (t^2 - 1)^n$：
\[
(2^nn!)^2\norm{P_n}^2 = \int_{-1}^1(u^{(n)})^2
= (-1)^n\int_{-1}^1 u\,u^{(2n)}
= (2n)!\int_{-1}^1(1 - t^2)^n\dd t ,
\]
（$u^{(2n)} = (2n)!$；边界项同第 4 题一样消失）。
又 $\int_{-1}^1(1-t^2)^n\dd t = B(\tfrac12, n+1) =
\frac{\Gamma(\frac12)\Gamma(n+1)}{\Gamma(n + \frac32)} =
\frac{2\cdot4^n(n!)^2}{(2n+1)!}$
（见\cref{exo:b3:product:8}）。综合得 $\norm{P_n}^2 =
\frac{2}{2n + 1}$.

\textbf{6.} 多项式在 $\norm\cdot_\infty$ 意义下
$\mathcal C(\intcc{-1}1)$ 中稠密（魏尔斯特拉斯，
\cref{cor:b3:complete:weierstrass}），连续函数在
$L^2$ 意义下稠密（\cref{thm:b3:lp:density}），且
$\norm\cdot_2 \leq \sqrt2\norm\cdot_\infty$：多项式张成
的子空间稠密，故归一化的 $P_n$ 构成\omterm{def:b3:hilbert:onb}{希尔伯特基}。
$\abs t$ 的展开中，与 $P_0$ 对应的系数为
$\frac{\langle P_0, \abs t\rangle}{\norm{P_0}^2} = \frac12$;
与 $P_1$：$0$（由奇偶性）；与 $P_2$：
$\frac{\int_{-1}^1\abs t\,\frac{3t^2-1}2\dd t}{2/5} =
\frac{1/4}{2/5} = \frac58$。最佳二次逼近为：
\[
\abs t \approx \frac12 + \frac58\,P_2(t) = \frac{3}{16} +
\frac{15}{16}\,t^2 .
\]

\textbf{7.} 由 $\frac{\dd^{n+1}}{\dd t^{n+1}}\eu^{-t^2} =
\frac{\dd^n}{\dd t^n}(-2t\,\eu^{-t^2})$ 及莱布尼茨公式，$H_{n+1}
= 2tH_n - H_n'$，由归纳法得次数为 $n$，首项
系数为 $2^n$。对 $m < n$，分部积分 $n$ 次，
即 $\int H_m H_n\eu^{-t^2} = (-1)^n\int H_m\,\bigl(\eu^{-t^2}
\bigr)^{(n)}$：边界项（多项式 $\times$
$\eu^{-t^2}$）在 $\pm\infty$ 处为零，剩下 $\int
H_m^{(n)}\,\eu^{-t^2} = 0$。对 $m = n$：$H_n^{(n)} = 2^nn!$，
故 $\norm{H_n}_w^2 = 2^nn!\int\eu^{-t^2} = 2^nn!\sqrt\pi$。

\textbf{8.} 设 $f \in L^2(\R, \eu^{-t^2}\dd t)$ 与
每个多项式正交，令 $g = f\eu^{-t^2}$。则 $g \in
L^1$：$\int\abs f\eu^{-t^2} \leq \bigl(\int\abs
f^2\eu^{-t^2}\bigr)^{1/2}\bigl(\int\eu^{-t^2}\bigr)^{1/2}$
（柯西--施瓦茨不等式）。对 $\xi \in \R$，展开 $\eu^{-\iu\xi t}$：
部分和被控制，因为
\[
\sum_k\frac{\abs\xi^k}{k!}\int\abs f\,\abs t^k\eu^{-t^2}\dd t
\leq \Bigl(\int \abs f^2\eu^{-t^2}\Bigr)^{1/2}
\sum_k\frac{\abs\xi^k}{k!}\Bigl(\int
t^{2k}\eu^{-t^2}\Bigr)^{1/2} < \infty
\]
（最后的级数收敛：$\int t^{2k}\eu^{-t^2} =
\Gamma(k+\frac12) \leq k!\,\sqrt\pi$，故各项为
$O(\abs\xi^k/\sqrt{k!})$）。逐项积分
（对绝对值级数应用\cref{cor:b3:lebesgue:additivity}，
再对级数用富比尼定理）得
\[
\int_\R g(t)\,\eu^{-\iu\xi t}\dd t
= \sum_k\frac{(-\iu\xi)^k}{k!}\int f(t)\,t^k\,\eu^{-t^2}\dd t
= 0 ,
\]
每个积分都是 $\langle t^k, f\rangle_w$ 型，$= 0$。由
在本阶段不加证明地承认的傅里叶变换单射性
（见\cref{ch:b3:fouriertransform}），$g = 0$ 几乎处处，故 $f = 0$
几乎处处，故埃尔米特族（其张成为全体多项式）是
稠密的。

\textbf{9.} 精确至次数 $n - 1$：对这样的 $P$，$P =
\sum_iP(t_i)\ell_i$ 精确成立，故 $\int Pw = \sum_iP(t_i)\int
\ell_iw = Q(P)$。次数 $\leq 2n - 1$ 时，作带余除法 $P = qp_n + r$，
$\deg q \leq n - 1$，$\deg r \leq n-1$；于是 $\int Pw = \int
qp_nw + \int rw = 0 + Q(r)$（$p_n \perp \R_{n-1}[t]$），而
$Q(P) = \sum_iw_i\bigl(q(t_i)\,p_n(t_i) + r(t_i)\bigr) = Q(r)$
因节点为 $p_n$ 的根，两者相等。

\textbf{10.} $\ell_i^2$ 的次数为 $2n - 2 \leq 2n - 1$，且
$\ell_i^2(t_j) = \delta_{ij}$: $0 < \int\ell_i^2w = Q(\ell_i^2)
= w_i$。由波利亚准则（见\cref{exo:b3:banach:9}，移至区间 $I$ 上并带权
）：条件 (i) 成立——每个多项式都被
精确积分一次（只需 $2n - 1 \geq$ 其次数）；条件 (ii)：
$\sum_i\abs{w_{i}} = \sum_iw_i = Q(\mathbf 1) = \int_Iw$,
有界：$Q_n(f) \to \int fw$ 对每个 $f \in \mathcal C(I)$，
$I$ 紧。

\textbf{11.} 首一多项式 $p_2 = t^2 - \frac13$（由
\cref{exo:b3:hilbert:4}）：节点为 $\pm\frac1{\sqrt3}$。权：
$\ell_1(t) = \frac{t - \frac1{\sqrt3}}{-\frac2{\sqrt3}}$，且
$w_1 = \int_{-1}^1\ell_1 = 1$；由对称性得 $w_2 = 1$。
精确性：$\int 1 = 2 = 1 + 1$；$\int t = 0 =
-\frac1{\sqrt3} + \frac1{\sqrt3}$; $\int t^2 = \frac23 =
\frac13 + \frac13$；$\int t^3 = 0$。两点梯形公式
（节点 $\pm1$，权 $1, 1$）仅精确至次数
$1$：对 $t^2$ 给出 $2$ 而非 $\frac23$。代价相同，
却多出两个精确次数：这正是正交
节点的收益。

\textbf{12.} 由 $\cos(n{+}1)\theta + \cos(n{-}1)\theta =
2\cos\theta\cos n\theta$：$T_{n+1} = 2tT_n - T_{n-1}$，其中
$T_0 = 1$，$T_1 = t$；由归纳法可得次数为
$n$ 的多项式，首项系数为 $2^{n-1}$（$n \geq 1$）。
作代换 $t = \cos\theta$（$w(t)\dd t \mapsto
\dd\theta$）：$\langle T_m, T_n\rangle_w =
\int_0^\pi\cos m\theta\cos n\theta\,\dd\theta = 0$，当 $m
\neq n$ 时，$= \pi$ 对 $m = n = 0$，$= \frac\pi2$ 其余情形
（积化和差）。由次数及两两正交性，
可知 $T_n$ 与格拉姆--施密特正交化的输出
仅相差数量因子；$f$ 的切比雪夫展开恰是
$\theta \mapsto f(\cos\theta)$ 的傅里叶余弦级数。

\textbf{13.} $T_n(t) = 0$ 当且仅当 $\cos n\theta = 0$，
亦即 $\theta = \frac{(2k-1)\pi}{2n}$：$n$ 个互异的根
$t_k = \cos\frac{(2k-1)\pi}{2n} \in \intoo{-1}1$。极值：
$\abs{T_n} \leq 1$ 在 $\intcc{-1}1$ 上，且 $T_n(s_j) =
(-1)^j$，在 $n + 1$ 个点 $s_j = \cos\frac{j\pi}n$ 处取到：
这是完美的等振荡。

\textbf{14.} $2^{1-n}T_n$ 为首一多项式，上确界范数为 $2^{1-n}$。
若某首一多项式 $P$ 的次数为 $n$，且 $\sup\abs P < 2^{1-n}$，则
差 $D = 2^{1-n}T_n - P$ 的次数至多为 $\leq n -
1$（首项相消），却在 $s_0 >
\dots > s_n$（此处 $2^{1-n}T_n = \pm2^{1-n}$ 压过
$P$）：至少 $n$ 个零点——$D \equiv 0$，矛盾。
等号情形的唯一性：同一 $D$ 满足
$(-1)^jD(s_j) \geq 0$；次数至多为 $\leq
n-1$ 的非零多项式不可能满足 $n$ 个弱交错的极值
条件而不具有 $n$ 个计重根（若
$D(s_j) = 0$，对某内点 $s_j$，则该零点应计为二重
（计数时因 $D$ 局部保持符号）：再次 $D
\equiv 0$.

\textbf{15.} 拉格朗日误差公式（罗尔定理，本科二年级）
给出 $f - L_nf = \frac{f^{(n)}(\xi_t)}{n!}\,\omega(t)$，故
一致误差至多为 $\frac{\norm{f^{(n)}}_\infty}
{n!}\,\sup\abs\omega$，且 $\omega$ 为 $n$ 次首一多项式：
由第 14 题，$\sup_{\intcc{-1}1}\abs\omega \geq 2^{1-n}$，
等号成立当且仅当节点为切比雪夫根。于是
最优界为 $\norm{f - L_nf}_\infty \leq
\frac{\norm{f^{(n)}}_\infty}{2^{n-1}n!}$。若取等距节点，
$\sup\abs\omega$ 在靠近
端点处指数级增大，且即便对 $\frac1{1 + 25t^2}$ 作插值，
当 $n \to \infty$ 时在那里发散（龙格现象）；
切比雪夫节点正是其解决之道。

\textbf{16.} 对 $T_n(\cos\theta) = \cos
n\theta$ 求导，得 $T_n'(\cos\theta) = \frac{n\sin n\theta}
{\sin\theta}$，其极限为 $n^2$（$\theta \to 0$ 时），
而为 $(-1)^{n+1}n^2$（$\theta \to \pi$ 时）：$\abs{T_n'(\pm1)} =
n^2$。内点处，$\abs{T_n'(t)} \leq
\frac{n}{\sqrt{1 - t^2}} = O(n)$：二次的爆破
只出现在端点处（伯恩斯坦内点界相对于
马尔可夫整体界）。

\textbf{17.} 设 $\theta_k = \frac{(2k-1)\pi}{2n}$，且 $S_j
= \sum_{k=1}^n\cos(j\theta_k)$，$1 \leq j \leq n-1$。则
\[
S_j = \operatorname{Re}\Bigl[\eu^{\iu j\pi/2n}
\sum_{k=0}^{n-1}\eu^{\iu jk\pi/n}\Bigr]
= \operatorname{Re}\Bigl[\eu^{\iu j\pi/2n}\,
\frac{\eu^{\iu j\pi} - 1}{\eu^{\iu j\pi/n} - 1}\Bigr] .
\]
当 $j$ 为偶数时分子为零：$S_j = 0$。当 $j$ 为奇数时
分子为 $-2$，且 $\eu^{\iu j\pi/n} - 1 =
\eu^{\iu j\pi/2n}\cdot2\iu\sin\frac{j\pi}{2n}$，于是整个
表达式为 $\frac{-2}{2\iu\sin(j\pi/2n)} =
\frac{\iu}{\sin(j\pi/2n)}$：纯虚数，$S_j = 0$
同样成立。由此，等权公式 $\frac\pi n\sum_kf(t_k)$
对 $T_0$ 精确（$\sum w_i = \pi = \int w$），并使 $T_1,
\dots, T_{n-1}$ 的值与 $\int T_jw = 0$ 一样为零，即精确至次数
$n - 1$。给定节点时，精确至次数 $n-1$ 的权
唯一（拉格朗日基）：高斯权全部为
$\frac\pi n$。当 $n = 3$ 时，节点为 $\pm\frac{\sqrt3}2, 0$，并有
\[
\int_{-1}^1\frac{f(t)}{\sqrt{1 - t^2}}\,\dd t \approx
\frac\pi3\Bigl[f\Bigl(\tfrac{\sqrt3}2\Bigr) + f(0) +
f\Bigl(-\tfrac{\sqrt3}2\Bigr)\Bigr],
\]
精确至次数 $5$。

\textbf{18.} 设 $P$ 为 $n$ 次首一多项式：$P = p_n + r$
其中 $r \in \R_{n-1}[t]$，且 $p_n \perp \R_{n-1}[t]$
（第 1 题），故 $\norm P^2 = \norm{p_n}^2 + \norm r^2 \geq
\norm{p_n}^2$，等号成立当且仅当 $r = 0$：$p_n$ 是
$t^n$ 在
$\R_{n-1}[t]^\perp$ 上的正交投影余项，即与
其必须避开的子空间最近的首一多项式。切比雪夫的 $2^{1-n}T_n$ 回答了
同一问题（对上确界范数）：与零的偏差最小，
一次在 $L^2(w)$ 中，一次在 $L^\infty$ 中。

\textbf{19.} 记 $K_n(x, y) =
\sum_{k=0}^n\frac{p_k(x)p_k(y)}{h_k}$。基础情形 $n = 0$：$(x -
y)\frac1{h_0} = \frac{p_1(x)\cdot1 - 1\cdot p_1(y)}{h_0}$
因 $p_1 = t - a_0$。归纳步：设恒等式对 $n -
1$ 成立，
\[
(x - y)\,K_n(x,y) = \frac{p_n(x)p_{n-1}(y) -
p_{n-1}(x)p_n(y)}{h_{n-1}} +
\frac{(x - y)\,p_n(x)p_n(y)}{h_n} ;
\]
在第二项中代入 $x\,p_n(x) = p_{n+1}(x) + a_np_n(x) +
b_np_{n-1}(x)$ 及 $y\,p_n(y) = p_{n+1}(y) + a_np_n(y) +
b_np_{n-1}(y)$：$a_n$ 的贡献
相消，而 $b_n = \frac{h_n}{h_{n-1}}$ 的贡献
与归纳项相消；剩下的是
$\frac{p_{n+1}(x)p_n(y) - p_n(x)p_{n+1}(y)}{h_n}$。
合流形式由令 $y \to x$ 得出（两边均为
$y$ 的多项式）。

\textbf{20.} 由合流形式得 $p_{n+1}'p_n -
p_n'p_{n+1} = h_n\sum_{k\leq n}\frac{p_k^2}{h_k} \geq
\frac{h_n}{h_0} > 0$，处处成立。设 $x_0$ 为
$p_{n+1}$ 的根，则 $p_{n+1}'(x_0)\,p_n(x_0) > 0$，故 $p_n(x_0) \neq
0$（无公共根）。设 $x_0 < x_1$ 为
$p_{n+1}$ 的相邻两根（均为单根，见第一部分），$p_{n+1}'$ 在两根处
符号相反，从而 $p_n$ 亦然：$p_n$ 的一根落在
这 $n$ 个间隔中的每一个——而这恰好用尽了它的 $n$ 个根：
零点交错。

\textbf{21.} 将 $D_n(t) = \det(tI_n - J_n)$ 沿
最后一行展开：$D_n = (t - a_{n-1})D_{n-1} - b_{n-1}D_{n-2}$，
其中 $D_0 = 1$，$D_1 = t - a_0$：这正是
首一多项式 $p_n$ 的递推关系与初值，故 $D_n = p_n$。$p_n$ 的根即
对称矩阵 $J_n$ 的特征值：实数，且由
第 19 题知其为单根——高斯求积实为
伪装成三对角矩阵的谱理论，是
在\cref{ch:b3:spectral}中谱理论的有限维影子。

\textbf{22.} 对照表：
\begin{center}
\begin{tabular}{l|c|c|c}
 & 勒让德 & 埃尔米特 & 切比雪夫 \\
\hline
区间 & $\intcc{-1}1$ & $\R$ & $\intcc{-1}1$\\
权函数 & $1$ & $\eu^{-t^2}$ & $(1-t^2)^{-1/2}$\\
公式 & 罗德里格斯 & $(-1)^n\eu^{t^2}
\frac{\dd^n}{\dd t^n}\eu^{-t^2}$ & $\cos(n\arccos t)$\\
范数$^2$ & $\frac2{2n+1}$ & $2^nn!\sqrt\pi$ & $\pi,
\frac\pi2$\\
适用领域 & 求积 & 高斯积分 & 极小极大 \\
\end{tabular}
\end{center}
（各族均有其三项递推式，勒让德族的一般形式为
$H_{n+1} = 2tH_n - 2nH_{n-1}$，$T_{n+1} = 2tT_n -
T_{n-1}$）。一般理论提供了任何单一族都
无法显示的内容：根的实性与交错性，求积
权的正性，递推关系的存在，以及
克里斯托费尔--达布公式——这些都仅由正交性推出，
与权函数无关。

\textbf{23.} 存在性：线性映射 $\R_{2n-1}[t] \to
\R^{2n}$，$P \mapsto (P(t_1), P'(t_1), \dots, P(t_n),
P'(t_n))$ 为单射（核中的 $P$ 有 $n$ 个二重
根且次数 $\leq 2n - 1$，故 $P = 0$），两空间之间
维数相同，均为 $2n$，故为双射。逐点误差：固定非节点的 $t$，
并选取 $K$，使 $g(s) = f(s) - Hf(s) -
K\,p_n(s)^2$ 在 $s = t$ 处为零。于是 $g$ 在
$n + 1$ 个互异点 $t, t_1, \dots, t_n$ 上为零，且 $g'$
在每个 $t_i$ 处亦为零（$f - Hf$ 与 $p_n^2$ 在那里都有
二重零点）。由罗尔定理得 $n$ 个 $g'$ 的零点，严格
位于 $g$ 的相邻零点之间——且异于节点
——因而 $g'$ 共有 $2n$ 个互异零点；再反复应用罗尔定理 $2n - 1$
次，得到 $\xi_t$，使 $g^{(2n)}(\xi_t) = 0$。
由于 $\deg Hf \leq 2n - 1$，且 $p_n^2$ 为首一多项式，次数为
$2n$，故 $g^{(2n)} = f^{(2n)} - K\,(2n)!$，于是 $K =
f^{(2n)}(\xi_t)/(2n)!$——恒等式在
节点处显然成立。积分：$Q_n(f) = Q_n(Hf)$（$Hf$ 与 $f$ 在
节点处一致），且 $Q_n(Hf) = \int Hf\,w$，由精确至次数
$2n - 1$（第 9 题），故求积误差为
$\int(f - Hf)\,w$。记 $m, M$ 为 $f^{(2n)}$ 在
$I$ 上的最小值与最大值，由逐点恒等式得如下夹逼：
\[
\frac{m\,h_n}{(2n)!} \;\leq\; \int_I(f - Hf)\,w
\;\leq\; \frac{M\,h_n}{(2n)!} ,
\]
再由介值定理，对连续函数
$f^{(2n)}$ 得出 $\xi$。（勒让德情形，$n = 2$：$h_2
= \int_{-1}^1(t^2 - \frac13)^2\dd t = \frac8{45}$，故
误差为 $f^{(4)}(\xi)/135$。）

\textbf{24.} 该核对 $\R_{n-1}[t]$ 具有再生性：将
$q = \sum_k\frac{\langle p_k, q\rangle}{h_k}p_k$ 展开，得
$\int_I K_n(t_i, t)\,q(t)\,w(t)\dd t = q(t_i)$ 对每个
$q$（次数 $\leq n - 1$）成立。取 $q = \ell_i$，则左边
等于 $\ell_i(t_i) = 1$。但 $t \mapsto K_n(t_i,
t)\,\ell_i(t)$ 为次数 $\leq (n - 1) + (n
- 1) = 2n - 2$ 的多项式，$Q_n$ 对其精确（第 9 题），且它
在每个节点 $t_j \neq t_i$ 处为零（因子 $\ell_i$），故
\[
1 = \int_I K_n(t_i, t)\,\ell_i(t)\,w(t)\dd t
= w_i\,K_n(t_i, t_i)
= w_i\sum_{k=0}^{n-1}\frac{p_k(t_i)^2}{h_k} .
\]
和式 $> 0$（其 $k = 0$ 项为 $1/h_0 > 0$）：即所述的
公式，并再次说明正性。验证（$n = 2$，
勒让德情形）：$p_0 = 1$，$h_0 = 2$，$p_1 = t$，$h_1 = \frac23$；
在 $t_i = \pm\frac1{\sqrt3}$ 处，
\[
K_2(t_i, t_i) = \frac12 + \frac{1/3}{2/3} = 1,
\qquad w_i = 1,
\]
如第 11 题所得。

\textbf{25.} 作代换 $t = \cos\theta$，积分为
$\int_0^\pi\cos^6\theta\,\dd\theta =
\pi\,\frac{5\cdot3\cdot1}{6\cdot4\cdot2} = \frac{5\pi}{16}$
（沃利斯，见\cref{exo:b3:product:8}）。$n = 3$ 的
切比雪夫--高斯求积公式（第 17 题）的节点为
$\cos\frac\pi6 = \frac{\sqrt3}2$，$\cos\frac\pi2 = 0$，
$\cos\frac{5\pi}6 = -\frac{\sqrt3}2$，且权相等，均为
$\frac\pi3$：
\[
Q_3(t^6) = \frac\pi3\Bigl(2\cdot\Bigl(\frac{\sqrt3}2
\Bigr)^{6}\Bigr) = \frac\pi3\cdot\frac{54}{64}
= \frac{9\pi}{32},
\qquad
\frac{5\pi}{16} - \frac{9\pi}{32} = \frac\pi{32} .
\]
预测：首一的 $3$ 次\omterm{pb:b3:hilbert:1}{正交多项式}为
$2^{-2}T_3 = t^3 - \frac34t$，其中 $h_3 =
\frac1{16}\norm{T_3}_w^2 = \frac1{16}\cdot\frac\pi2 =
\frac\pi{32}$；又 $f = t^6$，其 $f^{(6)} = 720 =
6!$ 为常数，故由第 23 题得误差为 $\frac{6!}{6!}\,h_3 =
\frac\pi{32}$——不再依赖 $\xi$，
该公式必然精确，事实上确实精确。
\end{solution}
